% Generated draft: requires mathematical and source review.
\documentclass[11pt]{article}
\usepackage[margin=25mm]{geometry}
\usepackage{fontspec}
\setmainfont{Latin Modern Roman}
\usepackage{amsmath,amssymb,mathtools,unicode-math}
\setmathfont{Latin Modern Math}
\usepackage{enumitem,needspace,xurl,hyperref}
\setlength{\parindent}{0pt}
\setlength{\parskip}{5pt}
\setlength{\emergencystretch}{3em}
\newcommand{\N}{\mathbb{N}}
\newcommand{\Z}{\mathbb{Z}}
\newcommand{\Q}{\mathbb{Q}}
\newcommand{\R}{\mathbb{R}}
\newcommand{\C}{\mathbb{C}}
\newcommand{\F}{\mathbb{F}}
\newcommand{\A}{\mathbb{A}}
\newcommand{\PP}{\mathbb{P}}
\newcommand{\E}{\mathbb{E}}
\newcommand{\eps}{\varepsilon}
\newcommand{\dd}{\,\mathrm{d}}
\begin{document}
\section*{UnsolvedMath ID: 30005454}
\textbf{Source problem code:} catalogue entry
\subsection{Definitions and mathematical statement}
Let \(G\) be the infinite path with vertex set \(\Z\) and edge set \(\{\{n,n+1\}:n\in\Z\}\). Set the reinforcement parameter to \(\alpha=1\), and let the firing rates satisfy \(p(v)=1\) for every \(v\in\Z\).

The supplied source denotes the relevant stochastic coordinates by \(x_v(t)\). Its Conjecture 2 asks whether
\[
\mathbb P\!\left(\bigcap_{v\in\Z}\left\{\lim_{t\to\infty}x_v(t)=1\right\}\right)=1.
\]
Thus convergence is required at every vertex on a single event of probability one; no uniform convergence over vertices is asserted. Since \(\Z\) is countable, this is equivalent to almost-sure convergence for each fixed vertex, once the stochastic process and its probability law have been specified.

\textbf{Unresolved specification.} The supplied record does not define the stochastic evolution, the normalization represented by \(x_v(t)\), the allowed initial configurations, the probability law, or the time parameter. Moreover, the surrounding excerpt uses edge-indexed coordinates \(x_e\) and \(X_e(t)\), whereas the conjecture uses vertex-indexed \(x_v(t)\). No identification between these coordinates is supplied. Consequently, the displayed assertion faithfully records the intended convergence target, but it is not a self-contained conjecture about a fully specified stochastic process. Its quantification over initial configurations must be recovered from the source rather than assumed.

The supplied motivation concerns the critical case \(\alpha=1\): a function \(L\), not defined in the record, is said to be non-strictly concave, permitting nonunique equilibria. With uniform firing rates, the quoted equilibria on a cycle of length \(2k\) have the alternating form \((a,b,a,b,\ldots,a,b)\), with \(a+b=2\); the admissible range of \(a,b\) is not specified. The excerpt states that odd cycles have a unique equilibrium and that, as the even cycle length increases, the law of \((a,b)\) converges to the Dirac measure at \((1,1)\). The law and initialization underlying this statement are not supplied. The interpretation of \(\Z\) as an ``infinite cycle'' is presented as motivation for the conjecture.
\subsection{Short English statement}
At critical reinforcement \(\alpha=1\), with every vertex of the infinite line firing at rate one, the supplied conjecture asks whether all coordinates \(x_v(t)\) converge to one almost surely. The process, normalization, initial conditions, and coordinate indexing require clarification before this becomes a self-contained mathematical statement. The record labels the problem open; that status is supplied curation, not independently verified here.
\subsection{Sources}
\begin{itemize}
\item Local UnsolvedMath registry: \texttt{data/unsolvedmath/problems.json}, permanent ID \(30005454\), exact source problem code \texttt{catalogue entry}, title ``Critical Reinforcement Convergence on the Infinite Line.'' The supplied record assigns proposed year 2023 and identifies no individual proposer.
\item Supplied source citation: \emph{MATRIX-MFO Tandem Workshop: Stochastic Reinforcement Processes and Graphs} (2023), Oberwolfach Reports, pp.\ 639--679, DOI \texttt{10.4171/owr/2023/12}. Original report: \url{https://doi.org/10.4171/owr/2023/12}. The supplied excerpt identifies the target as Conjecture 2 in the critical-case discussion.
\item Supplied literature reference: Couzinié--Hirsch, ``weakly reinforced Pólya urns,'' \url{https://arxiv.org/abs/2010.03347}. The surrounding excerpt attributes partial subcritical results to Y.\ Couzinié and C.\ Hirsch, reference [2], Theorem 2.2; the full numbered bibliography is absent.
\item Literature assessment dated \texttt{2026-08-22}: the supplied curation labels the problem open, states that the critical infinite-line conjecture appears in the report, and says that no proof of pointwise convergence to one on \(\Z\) was found in its review. It describes known results as concerning subcritical regimes or finite graphs. This is supplied dated curation, not an independently verified literature assessment.
\item The registry labels the statement \texttt{corrected\_verified} and reports a check against the official Oberwolfach report. This is a supplied source-verification claim; no independent source consultation or mathematical verification is asserted here.
\end{itemize}
\end{document}
